\begin{lemma}[A boundary edge can be adjoined exactly]
Let $G=(V,E)$ be a finite loopless multigraph, let $e_0\in E$, and let
$T$ be a tree-sequence with respect to $G$ and $e_0$.  If the exact edge
identity $f$ belongs to $\partial_G(V(T))$, then there is a tree-sequence
$U$ with respect to $G$ and $e_0$ such that
\[
  \operatorname{TreeSequenceAdjoinEdge}(T,f,U).
\]
In particular, $U$ has one more edge than $T$, contains $T$ as its ordered
initial segment, and has $f$ as its new final edge identity.
\end{lemma}
\begin{proof}
By the definition of the boundary in \noderef{BoundaryEdges}, choose ends
$x$ and $y$ of the exact edge identity $f$ such that
$x\in V(T)$, $y\notin V(T)$, and the unordered endpoint pair of $f$ is
$\{x,y\}$.  By \noderef{TreeSequenceVertexSet}, every displayed vertex
of $T$ belongs to $V(T)$.

First, $f$ is distinct from every displayed edge identity of $T$.  Indeed,
fix such an old edge $T.edge(i)$.  Its attachment clause in
\noderef{TreeSequence} says that its two ends are the newly displayed old
vertex at that step and an earlier displayed old vertex.  Both ends
therefore belong to $V(T)$.  If the exact identities $f$ and $T.edge(i)$
were equal, then $y$, being an end of $f$, would be one of those two ends
of $T.edge(i)$ and hence would belong to $V(T)$.  This contradicts the
choice of $y$.  This argument compares exact edge identities and does not
identify parallel edges merely because they have the same ends.

Write $p$ for the length of $T$.  Define an ordered sequence $U$ of length
$p+1$ by retaining, at their original indices, every vertex and edge of
$T$, and by putting $y$ at the new final vertex index and $f$ at the new
final edge index.  We verify that this data form a tree-sequence in the
sense of \noderef{TreeSequence}.  Two retained vertices are distinct by
vertex injectivity for $T$, while the new vertex $y$ is distinct from every
retained vertex because $y\notin V(T)$.  Thus the vertices of $U$ are
pairwise distinct.  Two retained edge identities are distinct by edge
injectivity for $T$, and the preceding paragraph shows that the new exact
edge identity $f$ is distinct from every retained edge.  Thus the edges of
$U$ are pairwise distinct.

The first edge of $U$ is the retained first edge of $T$, hence is $e_0$.
For every retained edge, retain its attachment witness from $T$ and regard
the same new-at-that-step and earlier vertices at their canonical old
indices in $U$; the endpoint equation is unchanged.  For the new final
edge, use $x\in V(T)$ and \noderef{TreeSequenceVertexSet} to choose an old
vertex index $q$ with $T.vertex(q)=x$.  The new final vertex is $y$, the
vertex at the retained index $q$ is $x$, and the endpoint pair of $f$ is
$\{x,y\}=\{y,x\}$.  Hence $f$ attaches the new final vertex to an earlier
displayed vertex.  Since $p\geq 1$ for the original tree-sequence, the new
length is positive.  All fields required of $U$ are now established.

It remains to check the claimed adjoining relation.  By construction,
$U$ has length $p+1$.  At every old vertex index and every old edge index,
the corresponding entry of $U$ is exactly the entry of $T$; hence $T$ is
an ordered initial segment of $U$ by
\noderef{TreeSequenceInitialSegment}.  At the unique new final edge index,
the edge of $U$ is exactly the identity $f$.  These are precisely the three
clauses of \noderef{TreeSequenceAdjoinEdge}, so
$\operatorname{TreeSequenceAdjoinEdge}(T,f,U)$ holds.
\end{proof}
